\refstepcounter{section}
\section*{Lecture 19: Schwarzschild solution}
\addcontentsline{toc}{section}{Lecture 19: Schwarzschild solution}
In the previous section, we calculated the various Christoffel symbols for the metric determining the orbit of planets around a star. We will continue further, towards a solution of the geodesic equation which will tell us about the motion.\\[0.2cm] As we had found out, the vacuum solutions require $\ST_{\mu\nu} = 0$ which implied that $R_{\mu\nu}=0$. Hence from the Christoffel symbol, we will now find the Ricci tensor and then solve the Einstein field equation. In general, this is an arduous (and boring) task to do it by hand. We will just mention the various non-zero components here \footnote{The appendix \ref{sec:appenB} contains code for symbolically computing the components of the Ricci and Riemann tensor} and carry on!
\begin{align}
    R_{tt} &= -\frac{1}{2h}\qty[f'' - \frac{f'h'}{2h} - \frac{(f')^2}{2f} + 2\frac{f'}{r}] \label{eqn:ric1} \\
    R_{rr} &= -\frac{1}{2f}\qty[f'' - \frac{(f')^2}{2f} - \frac{f'h'}{2h} - 2\frac{fh'}{rh}] \label{eqn:ric2} \\
    R_{\theta\theta} &= -\frac{r}{h^2f}\qty[hf - \qty(\frac{r}{h})^2 +1] \label{eqn:ric3}
\end{align}

Multiplying $2h$ with \eqref{eqn:ric1} and multiplying $2f$ with \eqref{eqn:ric2} and subtracting them, we get
\[
\frac{2h'f}{rh} + \frac{2f'}{r} = 0
\qquad \text{(since from field equation $R_{\mu\nu}=0$)}
\]

From this we get $$h'f+hf'=0\implies \dv{(hf)}{r} = 0 \longrightarrow h(r)f(r) = \SC_1 $$
where $\SC_1$ is a constant. Now we introduce some physical argument, that far away from the influence of the star, the metric should transform into the Minkowski metric (in polar coordinates obv) $$ds\tund{Mink.}^2 = -\ \dd t^2 +\ \dd r^2+ r^2\qty(\dd \theta^2 + \sin^2\theta \ \dd \phi^2)$$ since the effect of gravity becomes negligible as the potential falls off as $\sim r^{-2}$.
Thus we have:
% $$\lim\limits_{r\to \infty} h(r) = 1  \lim\limits_{r\to \infty} f(r) = -1 \begin{cases}
%      \lim\limits_{r\to \infty} h(r)f(r) = -1 \implies \lim\limits_{r\to\infty } \SC_1 = -1 \implies \SC_1 = -1
% \end{cases}$$


\begin{equation}\label{eqn:hf}
    \begin{rcases}
        \lim\limits_{r\to \infty} h(r) = 1 \\
        \lim\limits_{r\to \infty} f(r) = -1
    \end{rcases}
    \lim\limits_{r\to \infty} h(r)f(r) = -1 \implies \lim\limits_{r\to\infty } \SC_1 = -1 \implies \SC_1 = -1
\end{equation}
Combining \eqref{eqn:ric3} and \eqref{eqn:hf}, we get:
$$\dv{r}\qty(\frac{r}{h}) = 1 \longrightarrow \frac{r}{h} = r - r_s$$ 
where $r_s$ is an arbitrary constant which will be interpreted as a characterisitc radius (\textit{Schwarzschild radius}) later. Thus, we have essentially solved the problem upto an undetermined constant. 
$$h(r) = \frac{r}{r-r_s}\qquad f(r) = \frac{r_s-r}{r}$$
We now write the Christoffel symbols in terms of $r_s$ since we have an expression of $h(r)$ and $f(r)$ now.
\[
\begin{aligned}
\chris{t}{t}{r} &= \frac{r_s}{2r(r-r_s)} &
\chris{r}{t}{t} &= \frac{r_s(r-r_s)}{2r^3} &
\chris{r}{r}{r} &= \frac{r_s}{2r(r_s-r)} \\[2mm]
\chris{r}{\theta}{\theta} &= r_s-r &
\chris{r}{\phi}{\phi} &= (r_s-r)\sin^2\theta  &
\chris{\theta}{\phi}{\phi} &= -\sin\theta\cos\theta \\[1mm]
\chris{\phi}{\theta}{\phi} &= \cot\theta &
\chris{\phi}{\phi}{r} &= \frac{1}{r} &
&
\end{aligned}
\]

The motion can now be obtained from the geodesic equation: 
$$\dv[2]{{x^\mu}}{{\tau}} + \chris{\mu}{\alpha}{\beta}\dv{{x^\alpha}}{\tau}\dv{{x^\beta}}{\tau}=0 \qquad \mu \in \{t,r,\theta,\phi\}$$
$\vartriangleright$ \underline{Equation for Time}
\begin{alignat*}{3}
    &\dv[2]{{t}}{{\tau}} + \chris{t}{\alpha}{\beta}\dv{{x^\alpha}}{\tau}\dv{{x^\beta}}{\tau}&=0\\
    \implies\quad &\dv[2]{{t}}{{\tau}} + 2\frac{r_s}{2r(1-r_s)} \dv{{t}}{\tau}\dv{{r}}{\tau}&=0\\
    \implies\quad &\dv[2]{{t}}{{\tau}} + \frac{r_s}{r(1-r_s)}\dv{{t}}{\tau}\dv{{r}}{\tau}&= 0
\end{alignat*}
The above is the equation for time in the given metric. It feels weird since while constructing equation of motions, we had always written quantitites to vary with time but here, time itself varies with proper time. This equation essentially tells us how clock changes with respect to the proper time $\tau$. 